% Chapter 2
%\chapter{کار‌ها و نتایج پیشین}

%

\chapter{تعریف مفاهیم پایه}
\section{مقدمه}
در ابتدای مطالعه این پژوهش، آشنا شدن با تعاریف اولیه و مطالعه آن‌ها از زاویه دید‌های مختلف، می‌تواند به فهم پژوهش کمک کند. از این رو در این فصل، به این تعاریف پرداخته شده است.
\section{دوقلوی دیجیتال}

دوقلوی دیجیتال به عنوان یک فناوری نوین در حوزه مهندسی و مدیریت سیستم‌ها، از جمله تحولات مؤثر در عصر فناوری اطلاعات محسوب می‌شود. این مفهوم بیشتر در زمینه مدیریت چرخه عمر محصولات
\LTRfootnote{Product Lifecycle Management}
مورد توجه قرار گرفته است. 
در واقع، دوقلوهای دیجتیال به معنای وجود یک مدل مجازی وابسته به یک شیء فیزیکی در دنیای واقعی است که از طریق تبادل داده‌ها، همگام‌سازی و تجزیه و تحلیل آنچه در دنیای واقعی اتفاق می‌افتد، امکان بهینه‌سازی عملیات و پیش‌بینی مشکلات را فراهم می‌سازند. یک دوقلوی دیجیتال می‌تواند به روش‌های مختلفی با همتای حقیقی خود ارتباط برقرار کند و بر روی آن اثر گذاشته یا از آن اثر پذیرد.

دوقلوی دیجیتال به عنوان یک پیوند مؤثر میان دنیای فیزیکی و مجازی شناخته می‌شوند که نه تنها تبادل داده‌ را تسهیل می‌کند، بلکه بر پایه آن، امکان تعامل بلادرنگ و بهینه‌سازی فرآیندها را نیز ایجاد می‌کند.
\cite{b2}
این مفهوم از طریق ایجاد یک اتصال دائمی و دوره‌ای بین فضای واقعی و مجازی به وجود می‌آید که در آن، داده‌های دنیای واقعی به فضای مجازی منتقل می‌شوند و برعکس. 
از این رو، دوقلوی دیجیتال می‌تواند حاوی اطلاعات دقیقی از وضعیت فعلی و پیش‌بینی‌های مربوط به شرایط آینده یک سیستم باشد.

می‌توان دوقلوی دیجیتال را به سه بخش اصلی تقسیم کرد:
\cite{b2}
\begin{enumerate}
	\item 
	بخش اول، بخش مجازی که در یک فضای مجازی نیز قرار دارد.
	\item 
	بخش دوم، بخش واقعی و حقیقی که در یک فضای حقیقی و واقعی قرار دارد.
	\item 
	بخش سوم، جریان داده‌ها میان بخش واقعی و بخش حقیقی وجود دارد.
\end{enumerate}

در واقع دوقلوی دیجیتال به مجموعه‌ای از مدل‌‌ها و ماشین‌های فیزیکی یا مجازی گفته می‌شود که عملکرد و چرخه حیات یک موجودیت فیزیکی، را شبیه‌سازی\LTRfootnote{simulating}، تقلید\LTRfootnote{emulating} کرده یا آینه‌وار\LTRfootnote{mirroring} نمایش می‌دهند. البته این دوقلو تنها به موجودیت‌های فیزیکی خلاصه نمی‌شود و می‌توان در لایه‌های مختلفی دوقلو‌های دیجیتال را بررسی کرد:
\begin{itemize}
	\item 
	دوقلوی دیجیتال در سطح اشیاء که در آن سعی بر مجازی‌سازی و ایجاد دوقلوی مجازی از اشیاء و موجودیت‌های مختلف است.
	\item 
	دوقلوی دیجیتال در سطح فرایند که از مجموعه‌ای از کنش‌ و برهم‌ کنش موجودیت‌های مختلف ایجاد می‌شود.
	\item 
	دوقلوی دیجیتال سامانه که از مجموعه‌ای از فرایند‌ها ایجاد می‌شود و قابلیت تعریف و تفسیر سامانه‌ها را دارد.
\end{itemize}

علاوه‌بر انواع مختلفی از دوقلو‌های دیجیتال که به آن‌ها اشاره شد، می‌توان موارد دیگری را نیز در نظر گرفت که با ترکیب کردن آن‌ها، قدرت ایجاد دوقلو‌های دیجیتال بزرگ‌تر، جامع‌تر و کامل‌تری را به صورت قابل توجهی افزایش می‌دهد به طوری که بتوان دوقلو‌های دیجیتال یک انسان یا ویژگی‌های آن‌را نیز ایجاد نمود.

دوقلوهای دیجیتال تنها یک مدل یا شبیه‌سازی ساده نیست، بلکه یک مدل زنده،‌ هوشمند و در حال تکامل و توسعه است که به عنوان یک همتای مجازی برای موجودیت حقیقی ایجاد می‌شود.

\section{عامل}

به طور کلی برای 
\textit{عامل}
یک تعریف واحد و یکتا که در همه زمینه‌ها کاربرد داشته باشد، یافت نمی‌شود. این موضوع به این دلیل است که در زمینه‌های مختلف، می‌تواند ویژگی‌های مختلفی داشته باشد و همچنین اشکال فیزیکی و غیر فیزیکی متنوعی را می‌تواند شامل شود، از ربات‌ها گرفته تا شبکه‌های کامپیوتر و حتی عامل‌های نرم‌افزاری.
یکی از تعاریفی که در حوزه هوش‌مصنوعی ارائه شده به صورت زیر است:

\textit{
	یک برنامه نرم افزاری است که می تواند با محیط خود تعامل داشته باشد، داده ها را جمع آوری کند و از داده ها برای انجام وظایف خود و تصمیم‌گیری در جهت رسیدن به اهداف از پیش تعیین شده استفاده کند.}
\cite{b4}

در برخی‌ از منابع، ویژگی‌های زیر برای یک عامل در نظر گرفته شده است:
\cite{b5}


\begin{description}
	\item[$\bullet$ قرارگیری\footnotemark:] \footnotetext{Situatedness}
	قرارگیری به معنای تعامل یک عامل با محیط از طریق حسگرها و عملگرها است. محیطی که یک عامل در آن حضور دارد، بخش جدایی‌ناپذیری از طراحی آن به شمار می‌رود و ورودی‌ها به‌طور مستقیم به‌عنوان نتیجه‌ی تعاملات عامل با محیط دریافت می‌شوند. این عامل‌ها به صورت مستقیم یا غیر مستقیم بر محیط تأثیر می‌گذارند و تاثیر می‌پذیرند.
	\item[$\bullet$ خودمختاری \footnotemark:] \footnotetext{Autonomy}
	خودمختاری توانایی یک عامل را برای انتخاب عمل‌های خود به‌طور مستقل و بدون هیچ مداخله‌ای از سوی دیگر عامل‌ها یا انسان‌ها نشان می‌دهد. این ویژگی عامل را در برابر تأثیرات خارجی محافظت می‌کند و حالت‌های داخلی آن را از ناپایداری ناشی از اختلالات خارجی مجزا\LTRfootnote{Isolate} می‌سازد. بدین ترتیب، عامل قادر است به‌طور مؤثری در محیط خود عمل کند و تصمیم‌های بهتری اتخاذ نماید.
	\item[$\bullet$ توانایی استنتاج\footnotemark:] \footnotetext{Inferential Capability}
	توانایی استنتاج به کمترین توانایی عامل جهت نتیجه‌گیری و تصمیم‌گیری با استفاده از داده‌های موجود، گفته می‌شود. عامی می‌تواند با استنتاج مشاهده‌ها از طریق تعمیم اطلاعات، به میزانی از درک وضعیت‌های مختلف برسد و بر اساس آن بهترین تصمیم را بگیرد.
	\item[$\bullet$ واکنش‌دهی\footnotemark:] \footnotetext{Responsiveness}
	واکنش‌دهی به معنای توانایی درک شرایط محیط و پاسخ به آن به‌طور به‌موقع است. این ویژگی برای بسیاری از برنامه‌های بی‌درنگ\LTRfootnote{Realtime} اهمیت بالایی دارد، زیرا عامل باید قادر باشد بلافاصله به تغییرات محیطی واکنش نشان دهد و ضمن پایش مداوم شرایط، تصمیمات مناسبی را اتخاذ کند.
	\item[$\bullet$ پیش‌فعالی\footnotemark:] \footnotetext{Pro-activeness}
	
	پیش‌فعالی به معنای توانایی عامل در پاسخ به رفتارهای فرصت‌طلبانه و گرفتن تصمیمات بهینه و هدفمند است. عامل باید بتواند به صورت فعال به فرصت‌ها و چالش‌های جدید پاسخ دهد و در محیط پویا و متغیر به‌صورت موثرتر عمل نماید.
	\item[$\bullet$ رفتار اجتماعی \footnotemark:] \footnotetext{Social Behaviour}
	رفتار اجتماعی به قابلیت تعامل عامل با منابع خارجی در مواقع ضروری اشاره دارد تا به اهداف خاص دست یابد. علی‌رغم اینکه تصمیم‌گیری عامل باید آزادانه و بدون مداخله خارجی باشد، باید بتواند به‌طور مؤثر با دیگر عامل‌ها و نهادهای انسانی ارتباط برقرار کند و دانش خود را به اشتراک بگذارد. این امر به عامل کمک می‌کند تا از تجربیات دیگران بیاموزد و به حل مسائل مشترک کمک کند، که خود به بهبود همکاری و کارایی کلی سیستم‌های چندعاملی\LTRfootnote{Multi Agent Systems (MAS)} منجر می‌شود.
\end{description}
علاوه‌بر ویژگی‌هایی که بیان شد، عامل‌ها می‌توانند بسته به کارکرد خود، ویژگی‌های دیگری نیز داشته باشند. همچنین بنابر عملکری که هر عامل باید داشته باشد، ممکن است برخی از ویزگی‌های بیان شده را دارا نباشد. پس هر عامل را می‌توان بر اساس این ویژگی‌ها سنجید و بررسی کرد.





\section{سیستم‌های چند ‌عاملی}
یک سیستم چندعاملی از تعدادی عامل ایجاد می‌شود که به صورت مستقل یا گروهی، با همکاری یا به صورت رقابتی، در یک محیط برای رسیدن به یک هدف یا اهداف مشخص، تشکیل شده است. این عامل‌ها می‌توانند دانش خود را با هم به اشتراک بگذارند و به کمک یکدیگر نتیجه و تصمیم بگیرند و یا هرکدام به صورت مستقل عمل کرده و بر روی محیط تاثیر گذارند.

سیستم‌ چندعاملی به دلیل قابلیت‌های منحصربه‌فردی که ارائه می‌دهد، قابل استفاده در زمینه‌ها و کاربرد‌های مختلف است. در اینجا به چند مورد از قابلیت‌های این سیستم ‌پرداخته‌ شده است:
\cite{b5}

\begin{itemize}
	\item
	انجام محاسبات و عمل‌ها به صورت موازی و غیر همزمان باعث افزایش سرعت و کارایی سیستم می‌شود.
	\item 
	هر زمان که یک یا چند عامل از سامانه خارج شود و دچار مشکل شود، کار کل سامانه مختل نمی‌شود و به کار خود ادامه می‌دهد. به این صورت میزان اطمینان‌پذیری و پایداری سامانه افزایش می‌یابد.
	
	\item
	در صورتی که نیاز باشد می‌توان عامل‌های تازه به سامانه اضافه کرد که این مورد انعطاف‌پذیر و مقیاس‌پذیری سامانه را فراهم می‌کند.
	
	\item 
	هزینه مربوط به هر عامل می‌تواند کمتر از سامانه‌های مرکزی و غیر توزیع‌شده باشد.
	\item 
	توسعه عامل‌ها به دلیل کوچک‌تر بودن و مشخص بودن کار آن‌ها، راحت‌تر و ساده‌تر و با هزینه کم‌تر انجام‌پذیر است.
	
\end{itemize}

در سیستم‌‌های چندعاملی می‌توان مسئله‌ها و وظیفه اصلی و بزرگ را به مسئله‌ها و وظیفه‌های کوچک‌تر تقسیم کرد و انجام هرکدام از آن‌ها را به عامل‌های جداگانه سپرد. با این کار می‌توان کار‌های بزرگی که انجام آن‌ها توسط یک عامل بسیار سخت و پیچیده  و هزینه‌بر است را به راحتی انجام داد.
\cite{b6}

\section{هوش مصنوعی در دوقلو‌ی دیجیتال}

امروزه کمتر حوزه‌ای می‌توان یافت که در آن هوش مصنوعی دخالت نداشته‌ باشد. امروزه در زمینه‌های مختلف، استفاده از روش‌ها و الگوریتم‌های هوش مصنوعی کاربرد فراوانی یافته است. 

هوش مصنوعی در اوایل ظهور خود، با اینترنت اشیاء\LTRfootnote{Internet of things (IoT)}
پیوند خورد و محصولات تازه‌ای وارد بازار کرد. تعریف هوش مصنوعی اشیاء\LTRfootnote{Artificial Intelligence of things (AIoT)}
از ترکیب فناوری هوش مصنوعی و اینترنت اشیاء ایجاد شده‌است. در گذشته دوقلوی دیجیتال و هوش مصنوعی دو مفهوم کاملا مجزایی دانسته می‌شدند اما امروزه با گسترش و تنوع چشم‌گیر حسگر‌ها و همچنین پررنگ شدن اهمیت تحلیل داده‌ و پاسخگویی و اقدام هوشمندانه و سریع بر اساس داده، قابل جدا سازی نیستند. هوش مصنوعی به عنوان بخشی از لایه‌های دوقلوی دیجیتال کاملا پذیرفته شده و در حال استفاده است. این موضوع را می‌توان در مثال‌ها و نمونه‌هایی که از دوقلو‌ی دیجیتال آورده‌ خواهد شد به وضوح، مشاهده کرد.

دوقلوی‌ دیجیتال با هوش مصنوعی از داده‌ فرآیند‌ها برای پیش‌بینی اینکه فرآیند در آینده چگونه خواهد بود، استفاده می‌کند. این عملکرد نسبت به گذشته دوقلوی دیجیتال که تنها یک شبیه‌سازی و نمایش سه‌بعدی از فرایند بود، بسیار کاربردی‌تر و مهم‌تر است.

\section{جمع‌بندی}
در این فصل اصطلاحات و عبارت‌های اصلی پژوهش تعریف شدند. در واقع هرکدام از مفاهیم ذکر شده، خود به تنهایی شامل تعریف جامع‌ای می‌شوند اما نکته قابل توجه آن است که وقتی این مفاهیم در کنار هم قرار می‌گیرند، می‌توانند مفاهیم و توانایی تازه‌ای ایجاد کرده و نیاز‌های اساسی‌تری را پاسخ دهند. در همین راستا، فصل بعد، به کاربرد‌های مختلفی که از ترکیب‌ شدن این مفاهیم در چارچوب دوقلوی دیجیتال به دست می‌آیند، پرداخته است.




\chapter{کاربرد‌های دوقلوی دیجیتال}

\section{مقدمه}
آشنا شدن با کاربرد‌های دوقلوی دیجیتال، در انتخاب معماری مورد استفاده برای پیاده‌سازی دوقلوی دیجیتال، موثر است.
دوقلوی دیجیتال کاربرد‌های متنوع‌ای در زمینه‌های مختلف دارد. پزشکی و سلامت، ساخت و تولید، شهر هوشمند، وسایل نقلیه هوشمند، صنعت نفت و گاز و پتروشیمی، تنها گوشه‌ای از زمینه‌های مورد استفاده در دوقلوی دیجیتال است که در این فصل به تعدادی از این کاربرد‌ها اشاره شده و در نهایت به کار‌ها و نتایج پژوهش‌های پیشین در این زمینه پرداخته می‌شود.
\cite{b14}
\section{دوقلوی دیجیتال پزشکی‌ و سلامت}

با نگاهی گذرا به اطراف،‌ می‌توانیم نمونه‌های فراوانی از دستگاه‌ها و حسگر‌هایی را مشاهده کرد که در حال نظارت و بررسی وضعیت سلامت و فیزیولوژیکی انسان‌ها هستند. این دستگاه‌ها با پایش و بررسی دائمی یا موقت وضعیت جسمانی افراد، آن‌ را گزارش کرده و تحلیل می‌کند. برای نمونه می‌توان به ساعت‌های هوشمند، دستگاه‌های تصویربرداری پزشکی، نرم‌افزار‌‌های دیجیتال سلامت، دماسنج اشاره کرد. هرکدام از این دستگاه‌ها داده‌های فراوانی که مربوط به بخش‌های مختلف بدن و سلامت افراد هست تولید می‌کنند. این داده‌ها ممکن است با هم در ارتباط باشند و به بررسی و نظارت و پایش\LTRfootnote{Monitoring} آن‌ها به صورت یکجا نیاز باشد. همچنین علاوه بر مشاهده نیاز به تحلیل داده‌ها و ارسال اطلاعات مورد نیاز به پزشک متخصص است. این کار می‌تواند به کمک دوقلوی دیجیتال به همراه دریافت اطلاعات از دستگاه‌های اینترنت اشیاء انجام شود. استفاده از دوقلوی دیجیتال در اهدای عضو، آموزش جراحی و کاهش خطر سایر روش‌های پزشکی، رو به توسعه است.


\section{شهرهای هوشمند}

دوقلوی دیجیتال شهر یک مفهوم جدید است که مدلی از برنامه‌ریزی و ساخت و ساز شهری برای توسعه پایدار آینده ارائه می‌دهد که به طور موثری نوآوری‌ها در فناوری دیجیتال را با مکانیسم‌های عملیاتی شهری ترکیب و مسیری امکان‌پذیر برای ارتقاء شهری فراهم می‌کند. 
با نقشه‌برداری‌های دقیق از شهر، ارتباط واقعیت و دنیای مجازی با یکدیگر و بازخورد هوشمند نمونه‌ فیزیکی و دیجیتالی شهر، فعالیت‌ها و پدیده‌های شهری به صورت کارآمد‌تر و پربازده‌تر انجام می‌شود. همچنین خدمات روزمره برای شهروندان راحت‌تر و فراگیر‌تر می‌شود و به ایجاد یک محیط پایدار کمک می‌کند.

در حالی ک شهر‌ها در دنیا، به سرعت در مسیر رشد هستند، در این مسیر دچار چالش‌هایی نیز می‌شوند. رکود اقتصادی، افزایش آلودگی‌های هوا و زیست محیطی، تراکم ترافیک شهری، بلایای طبیعی مانند طوفان، سیل و زلزله، می‌توانند نمونه‌هایی از چالش‌هایی باشد که حتی شهر‌های مدرن نیز با آن‌ها دست‌و پنجه نرم می‌کنند. دوقلوی دیجیتال شهر از منظر‌های مختلف، می‌تواند یک راه‌حل بالقوه برای سیاست‌گذاران و تصمیم‌گیران در جهان باشد که به‌وسیله آن، در جهت بهبود مدیریت شهری و ایجاد یک بوم‌سازگان\LTRfootnote{Ecosystem} شهری مساعد برای صنایع و مردم حرکت کنند.

استفاده از دوقلوی دیجیتال در شهر‌ها می‌تواند میزان صرفه‌جویی هزینه‌ها را افزایش دهد و بازگشت سرمایه را به سرعت، با خود به همراه داشته‌ باشد. پس می‌توان به سادگی پیش‌بینی کرد که در آینده برای هر شهر، مدیریت بحران‌های موجود و تحلیل آن به وسیله دوقلوی دیجیتال، یک امر ضروری است.

\section{صنعت ساخت و تولید}
فناوری دوقلوی دیجیتال، به طور گسترده در صنعت تولید استفاده می‌شود. دوقلو‌ی دیجیتال در صنعت تولید، فرایند‌ها و رخداد‌ها، پدیده‌ها و اتفاق‌ها را واضح و شفاف می‌کند. صنعت تولید و ساخت‌وساز را می‌توان مهد دوقلوی دیجیتال دانست زیرا اولین استفاده‌ها و نیازمندی‌ها در راستای صنعت۴\LTRfootnote{Industry 4} برای دوقلوی دیجیتال، در این صنعت ایجاد و انجام شده‌ است.

دوقلوی دیجیتال در صنعت تولید، می‌تواند با پویش یکپارچه و بررسی دائمی اجزاء و آلات موجود در صنعت، از پیشامد‌های ناگوار پیشرو جلوگیری کند و باعث کاهش چشم‌گیر هزینه‌های مختلف شود. استفاده از فناوری‌های مربوط به اینترنت اشیاء در این صنعت بسیار رایج است. بنابر این رفتن به سمت دوقلوی دیجیتال در آن بیشتر از زمینه‌های دیگر مورد انتظار است. همچنین خودکار‌سازی در این صنعت از اهمیت به خصوصی برخوردار است که با استفاده از دوقلوی دیجیتال می‌تواند فرایند‌های خودکار‌سازی کامل‌تر و به‌صرفه تر شود.


\section{خودرو}
دوقلوی دیجیتال می‌تواند مدل مجازی یک وسیله نقلیه فیزیکی را ایجاد کند. دوقلوی دیجیتال وسایل نقلیه می‌تواند با جمع‌آوری داده‌های مختلف از سنسور‌های متنوع موجود در آن، یک کپی مجازی از آن را ایجاد کند. با داشتن دوقلوی دیجیتال خودرو، می‌توان پیش‌بینی‌های لازم، جهت اتخاذ بهترین تصمیم به صورت خودکار یا توسط راننده، انجام شود. 

در همین راستا می‌توان به خودرو‌های خودران نیز اشاره کرد. این خودرو‌ها برای اینکه به مرحله خودران بودن برسند نیاز به ساعت‌ها آموزش توسط داده‌های واقعی و شبیه‌سازی‌های موجود هستند. ایجاد یک دوقلوی دیجیتال برای خودرو و آموزش آن توسط داده‌هایی که خود، می‌تواند دوقلوی دیجیتال از سایر خودرو‌ها و جاده‌ها باشد، این عمل را ممکن می‌کند. با این شرایط، می‌توان خودرو‌های خودران را ابتدا در دنیای مجازی ایجاد و رشد داد و پس از آنکه آزمون خود را در دنیای مجازی پس دادند و تا حد امکان بدون خطا عمل کردند، در دنیای واقعی شروع به کار کنند.

دوقلوی دیجیتال در خودرو‌ها را می‌توان به سطح شهر و خیابان‌ها آورد. به این صورت که یک دوقلوی دیجیتال از خودرو‌های موجود در خیابان به طوری که با هم در ارتباط باشند، ایجاد کرد. به عبارتی دیگر، یک دوقلوی دیجیتال چندعاملی که در آن عامل‌ها، خودرو‌ها هستند که با هم تبادل اطلاعات می‌کنند و بر اساس آن،‌ تصمیم‌گیری کرده و عملی را انجام می‌دهند. مثلا می‌توانند فاصله مناسب را از یکدیگره به صورت ایمن، حفظ کنند و یا مسیر‌های دارای ترافیک را انتخاب نکنند و به صورتی هرکدام مسیر خود را برگزیند، که کمترین ترافیک در سطح شهر ایجاد شود.

می‌توان دوقلوی دیجیتال خودرو را برای هر کدام از قطعات مکانیکی خودرو به صورت دقیق ایجاد کرد. می‌دانیم که قطعات خودرو در کنار هم می‌توانند یک وسیله نقلیه کامل را ایجاد کنند. پس یک دوقلوی دیجیتال چندعاملی می‌توان برای هر خودرو داشت که هر عامل آن، یک قطعه مهم و حیاتی در آن، در نظر گرفته می‌شود و به صورت مستقل، بررسی و پایش شرایط لحظه‌ای آن مهم است و در صورت پیش‌بینی ایجاد مشکل یا خرابی در آن، می‌توان در زمان مناسب، آن را تعمیر یا تعویض کرد پیش از آنکه هزینه کلان یا غیر قابل جبرانی ایجاد کند.

\section{انسان دیجیتال}
ایجاد انسان مجازی یا دیجیتال و دوقلوی دیجیتال انسان\LTRfootnote{Human Digital Twin} 
یکی دیگر از مواردی است که در سال‌های اخیر مورد توجه قرار گرفته است. این نوع از انسان، نه تنها به صورت آواتارهایی هوشمند بجای منشی، در خدمت شرکت‌ها هستند بلکه با رشد هوش مصنوعی، تا حدی توانسته‌اند کاملا حسی مشابه حس ارتباط با انسان واقعی را القاء کنند.

انسان دیجیتال به طور گسترده به عنوان یک برنامه رایانه‌ای تعریف می‌شود که یک ظاهر مانند انسان را ارائه می‌دهد. این انسان مجازی، در واقع گرافیک رایانه‌ای، بینایی رایانه، گفتار هوشمند، پردازش زبان طبیعی و سایر فناوری‌ها را یکپارچه می‌کند. می‌توان از آن برای تولید محتوای دیجیتال و تعامل انسان و رایانه برای کمک به بهبود کارآیی تولید محتوا و تجربه کاربر استفاده کرد. در واقع دوقلوی دیجیتال انسان، یک انسان در دنیای مجازی و غیر فیزیکی است که به صورت یک محصول جامع، با ابزار‌های مختلف مانند یادگیری عمیق،‌ پردازش گفتار و نوشتار، توانایی تعامل با انسان و غیر انسان ایجاد می‌شود و کاربرد‌های فراوانی در زمینه‌های مختلف می‌تواند داشته باشد.

\section{جمعیت مجازی}
با ایجاد دوقلوی دیجیتال فرد به طوری که کاملا مانند او فکر کند، حرکت کند، تصمیم بگیرد و حتی واکنش‌های فیزیکی و شیمیایی بدن مشابه داشته باشد و همچنین جسم آن، به صورت مشابه با نمونه واقعی در برابر دارو‌ها و خوراک‌ها عمل کند، می‌توان به سمت ایجاد یک جامعه مجازی حرکت کرد.

بنابر این تکثیر دوقلوی دیجیتال بدن انسان می‌تواند در بخش‌های مختلف از اهمیت بالایی برخوردار باشد. برای مثال اگر فرایند درمان یک بیماری خاص در نظر گرفته شود، می‌توان فرایند درمان را در سه مرحله بررسی کرد:
\begin{enumerate}
	\item 
ساخت نسخه‌های نامحدود از مدل‌های شبکه از انسان و همه عوامل مولکولی به صورت کاملا مشابه با انسان و بیماری مربوطه.
	\item 
درمان محاسباتی در دوقلوی دیجیتال با هزاران دارو به منظور شناسایی بهترین دارو برای بیماری مورد نظر.
	\item 
درمان بیمار با داروی مشخص شده پس از تست‌های بسیار در نمونه‌های دیجیتالی.
\end{enumerate}

\section{صنایع هوافضا}
ایجاد دوقلوی دیجیتال برای صنایع هوافضا به ویژه هواپیما و اجزای پیچیده آن را می‌تواند یکی از پر اهمیت‌ترین کاربرد‌های دوقلوی دیجیتال دانست زیرا که هزینه‌های خرابی و تعمیر در این صنعت بسیار بالا است و این درحالی است که به کمک دوقلوی دیجیتال و ایجاد دوقلوی دیجیتال قطعات صنعت هوافضا، می‌توان بسیاری از خرابی‌ها و مشکل‌های آن‌ها را پیش‌بینی کرده و از خرابی آن‌ها جلوگیری کرد.

همچنین استفاده از دوقلوی دیجیتال به عنوان شبیه‌ساز پرواز برای هواپیما‌ها و سایر اقلام پرنده، می‌تواند در پیش‌بینی حالت‌های مختلف پرواز موثر باشد. همچنین می‌توان از ایجاد دوقلوی دیجیتال آن‌ها در آموزش خلبان‌ها استفاده شود و می‌توان شبیه‌سازی شرایط سخت و غیرممکن پیش از اینکه در واقعیت ایجاد شود، را ساخت.

\section{صنایع نفت، گاز و پتروشیمی}
در این صنعت یا به طور کلی، صنعت هیدروکربن، طیف وسیعی از فناوری‌های دیجیتال را برای بهبود بهره‌وری، کارآیی و ایمنی در نظر می‌گیرد و در عین حال، هزینه‌های سرمایه و عملیاتی، خطرات بهداشتی و زیست محیطی به میزان قابل توجهی کاهش می‌یابد. این منابع بسیار ارزشمند هستند و در جهت جلوگیری از هدر رفت آن در فرایند‌های مختلف تولیدی و انتقال،‌ باید اقدام شود و دوقلوی دیجیتال پایگاه‌های استخراج و راه‌‌های انتقال می‌تواند گزینه بسیار مناسبی برای این کار باشد.

همچنین قرارگیری ایستگاه‌های انتقال در مکان‌های دور از دسترس و سختی بررسی آن و هزینه بالای مربوط به نگهداری و مراقبت از سامانه‌های مورد نیاز در این صنعت، استفاده از دوقلوی دیجیتال را بیشتر مفید می‌سازد.

\section{کار‌ها و نتایج پیشین}


بر اساس مقاله 
\cite{b7}
دوقلو‌های دیجیتال معمولا از سه بخش ایجاد می‌شود. این سه بخش عبارتند از فضای واقعی دارای شئ واقعی، فضای مجازی دارای شئ مجازی و ارتباط میان شئ مجازی و واقعی و انتقال داده ‌میان آن‌ها.

طبق ادعایی که در مقاله‌های 
\cite{b8}
و
\cite{b9}
صورت گرفته‌است، دوقلو‌ی دیجیتال تنها یک شبیه‌سازی ساده نیست. بلکه یک مدل زنده، هوشمند و در حال تکامل است. دوقلوی دیجیتال، چرخه عمر همتای حقیقی خود را دنبال می‌کند تا بتواند آن را نظارت\LTRfootnote{Monitor}
و کنترل\LTRfootnote{Control}
کرده و همچنین کارکرد آن را بهبود بخشد. دوقلوی دیجیتال به صورت مستمر با فضا و محیط واقعی در ارتباط بوده و به صورت بی‌درنگ با آن، داده‌ ردوبدل می‌کند.


مقاله 
\cite{b10}
به بررسی دوقلو‌های دیجیتال در اینترنت وسایل‌ نقلیه\LTRfootnote{Internet of Vehicles (IoF)}
پرداخته است. در واقع این مقاله، یک سکو\LTRfootnote{Platform}
جهت به اشتراک گذاری داده، بین وسایل نقلیه موجود در سطح شهر را معرفی می‌کند. در این مقاله، یک دوقلوی دیجیتال، به صورت یک بازتاب از وسیله نقلیه در فضای ابری تعریف می‌شود که داده‌ها را با وسیله نقلیه در هر لحظه همگام می‌کند. داشتن این دوقلوی دیجیتال در بستر ابری، باعث می‌شود که بتوان با به اشتراک‌گذاری داده‌ها بدون محدودیت فاصله و پهنای باند، کارایی رانندگی را افزایش داده و تصمیمات دقیق‌تری بر اساس داده‌ها گرفته‌ شود.\\
در نهایت مقاله 
\cite{b11}
یک چارچوب\LTRfootnote{Framework}
به نام 
\lr{MADTwin}\LTRfootnote{Multi Agent Digital Twin}
ارائه می‌دهد که برای مدل‌سازی دوقلوی دیجیتال برپایه عامل‌\LTRfootnote{Agent}
مناسب است. از این چارچوب برای هوشمند‌سازی یک انبار‌ توسط ربات‌ها استفاده کرده تا بتوان به طور بهینه هر ربات را به سوی اهداف محلی خود و اهداف کل سیستم به صورت مشارکتی یا مستقل هدایت کند.

\section{جمع‌بندی}
همانطور که در این فصل معرفی شد، کاربرد‌های متنوعی از دوقلوی دیجیتال وجود دارد که بسیار موثر و مورد نیاز است. برای اینکه این کاربرد‌های متنوع، به صورت عملی پیاده‌سازی شوند، نیاز است تا یک معماری مناسب و موثر برای آن انتخاب شود. در فصل بعدی، یک معماری دوقلوی دیجیتال به کمک سیستم‌های چند‌ عاملی معرفی می‌شود که می‌تواند در کاربرد‌های مختلف، مورد استفاده قرار گیرد.
